\documentclass[10pt,a4paper]{amsart}
\usepackage[margin=19mm]{geometry}
\usepackage{amsmath,amssymb,mathtools}
\usepackage[scr=rsfs,cal=euler]{mathalfa}
\usepackage{xfrac}
\usepackage[unicode,hidelinks,bookmarks=false]{hyperref}
\newcommand{\ie}{i.e.\ }
\newcommand{\cf}{cf.\ }
\newcommand{\resp}{respectively\ }
\newcommand{\Subsection}{\subsection}
\providecommand{\nobreakdash}{\nobreak\hbox{-}}
\setlength{\emergencystretch}{2em}
\multlinegap0pt
\usepackage{amssymb}
\usepackage{enumerate}
\usepackage{bm}
\usepackage{algorithm}\usepackage{algpseudocode}
\newtheorem{proposition}{Proposition}
\newcommand{\argmin}[1]{\underset{#1}{\operatorname{arg}\operatorname{min}}\;}
\newcommand{\bx}{\mbf{x}}
\newcommand{\mbf}[1]{\boldsymbol{#1}}
\renewcommand{\top}{T}
\title{A Sparse Bayesian Learning Algorithm for Estimation of Interaction Kernels in Motsch-Tadmor Model}
\author{Jinchao Feng, Sui Tang}
\date{Source: arXiv:2505.07068v2 (CC BY 4.0)}
\begin{document}
\maketitle

\noindent\emph{Source and licence.} A Sparse Bayesian Learning Algorithm for Estimation of Interaction Kernels in Motsch-Tadmor Model. Version arXiv:2505.07068v2 (CC BY 4.0), distributed by arXiv under the Creative Commons Attribution 4.0 International licence, http://creativecommons.org/licenses/by/4.0/, as recorded in the arXiv licence metadata for this exact version and shown on the abstract page https://arxiv.org/abs/2505.07068v2. Full licence notice: \texttt{licenses/arxiv-2505.07068-CC-BY-4.0.txt}.\par\smallskip

\noindent\textit{Editorial scope.} Counted unit: the article's proposition prop1 (source0.tex lines 534--544). The article's complete original proof of it is delivered with it (source0.tex lines 546--556, one \texttt{proof} environment of 11 lines). No mathematics is written, repaired or re-derived: the statement and the proof are the article's own lines, in the article's own notation, and no retained line is substituted or re-flowed.  No further source-local context is needed: every cross-reference of every retained line resolves inside the two delivered spans.  Nothing was omitted from any retained span, so the package carries no omission record.  No retained line contains a citation, so the wrapper carries no bibliography and no bibliography entry.  No retained line contains a figure environment, an image-inclusion command or drawing code, so no image is delivered and the package declares no asset dependency.  Editorial formatting only, in addition: an AMS \texttt{amsart} preamble that restates the article's own declarations and macros used by the retained lines, namely the environments \texttt{proposition} on separate counters, together with the standard packages those lines need.  The retained environments print in this document as Proposition 1 in order of appearance, whereas the article prints the same items with the numbering of its own full document.  The complete native source of the article as distributed in the arXiv e-print for this exact version is bundled unchanged as \texttt{sources/177929/source0.tex}.
\begin{proposition}\label{prop1}
Let \(\hat{\mbf{c}}_{M,L} = \argmin{\|\mbf{c}\|_2=1} \mathcal{E}_{M,L}(\mbf{c})\). Then:
\begin{equation}\label{eq:prob1}
\frac{1}{NML} \mathbb{A}_{M,L}^{\top} \mathbb{A}_{M,L} \hat{\mbf{c}}_{M,L} = \mbf{0},
\end{equation}
where \(\mathbb{A}_{M,L} = \begin{pmatrix}\mbf{A}^{(1,1)} \\ \vdots \\ \mbf{A}^{(M,L)} \end{pmatrix} \in \mathbb{R}^{dNML \times K}\) with \(\mbf{A}^{(m,l)} = [\mbf{A}_{ik}^{(m,l)}] \in \mathbb{R}^{dN \times K}\), and
\begin{equation}\label{eq:assemble_A}
\mbf{A}_{ik}^{(m,l)} = \sum_{j=1}^{N} \xi_k(|\bx_j^{(m)}(t_l) - \bx_i^{(m)}(t_l)|) (\dot\bx_i^{(m)}(t_l) - (\bx_{i}^{(m)}(t_l) - \bx_{j}^{(m)}(t_l))) \in \mathbb{R}^d,
\end{equation}
for \(i=1,\cdots,N\) and \(k=1,\cdots,K\).
\end{proposition}

\begin{proof}
The loss function \(\mathcal{E}_{M,L}(\mbf{c})\) can be written as a quadratic form: 
\[
\mathcal{E}_{M,L}(\mbf{c}) = \frac{1}{NML} \|\mathbb{A}_{M,L} \mbf{c}\|^2 = \frac{1}{NML} \mbf{c}^\top \mathbb{A}_{M,L}^\top \mathbb{A}_{M,L} \mbf{c}.
\]
Differentiating with respect to \(\mbf{c}\) and setting it to zero gives:
\[
\mathbb{A}_{M,L}^\top \mathbb{A}_{M,L} \mbf{c} = \mbf{0},
\]
which proves the result.
\end{proof}

\end{document}
